% chapters/11/section-03.tex
\section{Batasan Kinerja dan Hitung Ulang Berkala}
\label{sec:batas-kinerja}

Pencarian jalur adalah perhitungan termahal pada blok Python ini. Bagian ini
menghitung besarnya biaya itu dan menjelaskan cara mengendalikannya.

\subsection{Menghitung Beban Perhitungan}

Biaya pencarian jalur bertambah menurut tiga hal: banyaknya NPC, seringnya
pencarian diulang, dan mahalnya satu kali pencarian. Ketiganya berlipat, bukan
menjumlah.

Sebagai gambaran, ambil angka dari pengukuran pada Bab~X, yaitu satu kali
pencarian memeriksa sekitar 90 sampai 280 petak. Dengan tiga NPC yang
menghitung ulang setiap bingkai pada laju 60 bingkai per detik, beban
perhitungannya adalah

\[
3 \times 60 \times 90 = 16.200 \text{ petak setiap detik},
\]

yaitu sekitar enam belas ribu kali pemeriksaan petak setiap detik. Angka itu
belum termasuk pekerjaan lain, yaitu menggambar, memperbarui gerak, dan
menangani tabrakan.

\begin{catatabox}
Perhitungan di atas memakai angka \textbf{per detik}, bukan
per bingkai. Kebiasaan berpikir dalam satuan per detik---yang telah dibangun
sejak Bab~IV---sangat berguna di sini, sebab beban perhitunganlah yang
menentukan apakah permainan masih dapat berjalan pada laju bingkai yang wajar.
\end{catatabox}

\subsection{Mengulang Pencarian Menurut Selang Waktu}

Cara paling langsung untuk menekan beban itu adalah memperlambat pengulangan
pencarian. Hasil pencarian disimpan pada NPC, lalu dipakai berkali-kali sampai
selang waktunya habis.

Pada program bab ini, selang itu dapat diatur melalui pilihan \code{--repath}.
Nilai nol berarti pencarian diulang setiap bingkai, sedangkan nilai yang lebih
besar memperlambat pengulangannya.

Memperlambat pengulangan bukan tanpa harga. Selama selang itu,
jalur yang ditempuh NPC menuju tempat pemain berada sebelumnya, sehingga NPC
sempat menempuh arah yang salah. Menentukan selang yang tepat karena itu
merupakan pertukaran, dan pertukaran itulah yang diukur pada
Bagian~\ref{sec:menilai-mutu}.

\subsection{Menyebar Waktu Perhitungan}

Satu persoalan lain muncul ketika jumlah NPC bertambah: apabila seluruh NPC
menghitung ulang jalurnya pada bingkai yang sama, beban perhitungan menumpuk
pada bingkai itu, lalu longgar pada bingkai-bingkai berikutnya. Gejalanya
adalah laju bingkai yang naik turun, walaupun rata-ratanya tetap.

Penanganannya sederhana: \textbf{berikan awal yang berbeda} kepada setiap NPC,
sehingga waktu perhitungan ulangnya tidak berbarengan.

\begin{kodekeluaran}
self.selang_tersisa = self.selang * (index / jumlah_npc)
\end{kodekeluaran}

Dengan pembagian itu, NPC pertama menghitung ulang pada awal selang, NPC kedua
pada sepertiga selang, dan seterusnya. Beban rata-ratanya sama, tetapi
sebarannya menjadi rata pula.

\begin{tipsbox}
Cara yang sama dipakai pada banyak bidang lain di luar permainan, misalnya
pada peladen yang membagi pekerjaan latar belakangnya, atau pada robot yang
membagi waktu perhitungan penglihatannya. Prinsipnya satu: pekerjaan yang
mahal tetapi tidak mendesak sebaiknya disebar, bukan ditumpuk.
\end{tipsbox}
